% ===========================================================================
%  Jàng Afrig — Lab 2 : HTML structurant
%  Sémantique, formulaires, navigation entre 5 pages
%  Licence 1 DUT1-Informatique — ESP/UCAD — 2025–2026
%  Dr. El Hadji Bassirou TOURÉ
%  Compilé avec : xelatex (2–3 passes)
% ===========================================================================
\documentclass[11pt,a4paper]{article}

\usepackage{fontspec}
\setmainfont{Latin Modern Roman}
\setsansfont{Latin Modern Sans}
\setmonofont{Latin Modern Mono}
\usepackage{polyglossia}
\setdefaultlanguage{french}

\usepackage[margin=2.5cm]{geometry}
\usepackage{amsmath,amssymb}
\usepackage{graphicx}
\graphicspath{{figures/}}
\usepackage{xcolor}
\usepackage{booktabs}
\usepackage{hyperref}
\usepackage{tcolorbox}
\tcbuselibrary{theorems,breakable,skins}
\usepackage{enumitem}
\usepackage{fancyhdr}
\usepackage{array}
\usepackage{pifont}
\usepackage{longtable}
\usepackage{multirow}

% --- TikZ pour les schémas ---
\usepackage{tikz}
\usetikzlibrary{shapes.geometric,arrows.meta,positioning,calc,backgrounds,fit,shadows.blur}

\newcommand{\cmark}{\ding{51}}
\newcommand{\xmark}{\ding{55}}

\hypersetup{
    colorlinks=true,
    linkcolor=blue!60!black,
    urlcolor=blue!60!black,
    citecolor=blue!60!black
}

\newtcolorbox{objectifbox}[1][]{colback=blue!5,colframe=blue!50,fonttitle=\bfseries,title=#1,breakable}
\newtcolorbox{infobox}[1][]{colback=green!5,colframe=green!50!black,fonttitle=\bfseries,title=#1,breakable}
\newtcolorbox{attentionbox}[1][]{colback=orange!5,colframe=orange!70,fonttitle=\bfseries,title=#1,breakable}
\newtcolorbox{retenirbox}[1][]{colback=yellow!10,colframe=yellow!50!black,fonttitle=\bfseries,title=#1,breakable}
\newtcolorbox{prereqbox}[1][]{colback=cyan!5,colframe=cyan!50!black,fonttitle=\bfseries,title=#1,breakable}
\newtcolorbox{projetbox}[1][]{colback=teal!5,colframe=teal!50!black,fonttitle=\bfseries,title=#1,breakable}

% --- Couleurs schémas ---
\definecolor{pageaccueil}{RGB}{45,130,200}
\definecolor{pagequiz}{RGB}{200,100,40}
\definecolor{pageresultat}{RGB}{200,80,60}
\definecolor{pagescores}{RGB}{120,80,180}
\definecolor{pageapropos}{RGB}{80,160,80}

% --- Code source HTML5 ---
\usepackage{listings}
\lstdefinelanguage{HTML5}{
  morekeywords={<!DOCTYPE,html,head,meta,title,body,h1,h2,h3,h4,p,ul,ol,li,table,tr,th,td,a,img,strong,em,br,div,span,header,nav,main,section,article,footer,aside,form,input,select,option,textarea,label,button,fieldset,legend},
  sensitive=true,
  morecomment=[s]{<!--}{-->},
  morestring=[b]",
}
\lstset{
  basicstyle=\ttfamily\small,
  backgroundcolor=\color{gray!8},
  frame=single,
  framerule=0.5pt,
  rulecolor=\color{gray!40},
  breaklines=true,
  breakatwhitespace=true,
  showstringspaces=false,
  tabsize=2,
  xleftmargin=4pt,
  xrightmargin=4pt,
  aboveskip=8pt,
  belowskip=8pt,
  literate={é}{{\'e}}1 {è}{{\`e}}1 {ê}{{\^e}}1 {à}{{\`a}}1 {ù}{{\`u}}1
           {ô}{{\^o}}1 {î}{{\^i}}1 {ç}{{\c{c}}}1 {É}{{\'E}}1
           {→}{{$\rightarrow$}}1
}

\pagestyle{fancy}
\fancyhf{}
\fancyhead[L]{\small Jàng Afrig --- Lab 2 : HTML structurant}
\fancyhead[R]{\small UCAD --- L1 DUT1-Info}
\fancyfoot[C]{\thepage}
\fancyfoot[L]{\small Dr. El Hadji Bassirou TOURÉ}
\fancyfoot[R]{\small 2025--2026}
\renewcommand{\headrulewidth}{0.4pt}
\renewcommand{\footrulewidth}{0.4pt}

\title{%
\vspace{-1cm}
\textbf{\LARGE Fondamentaux du Web}\\
\textbf{\LARGE HTML, CSS, JavaScript}\\[0.8cm]
\textbf{\huge Lab 2 --- HTML structurant}\\[0.3cm]
\textbf{\Large Sémantique, formulaires, navigation entre 5 pages}\\[0.5cm]
\large Licence 1 DUT1-Informatique --- 2025--2026}
\author{Dr. El Hadji Bassirou TOURÉ\\
Département Génie Informatique\\
Université Cheikh Anta Diop de Dakar}
\date{}

\begin{document}

\maketitle
\thispagestyle{empty}

\vspace{0.5cm}

% ===========================================================================
% BOÎTE PRÉREQUIS
% ===========================================================================
\begin{prereqbox}[Informations pratiques]
\begin{center}
\begin{tabular}{ll}
\toprule
\textbf{Type} & Lab guidé \\
\textbf{Séance} & S3 --- deuxième Lab du semestre \\
\textbf{Prérequis} & Lab 1 terminé (page d'accueil HTML poussée avec tag \texttt{v0.1}) \\
\textbf{Validation} & 5 pages HTML liées + formulaire fonctionnel + tag \texttt{v0.2} \\
\textbf{Rapport} & 1-2 pages : ce que j'ai fait + difficultés + lien repo \\
\bottomrule
\end{tabular}
\end{center}
\end{prereqbox}

\begin{objectifbox}[Objectifs du Lab]
À la fin de ce Lab, vous aurez :
\begin{enumerate}[nosep]
    \item Compris et utilisé les \textbf{balises sémantiques} HTML5 (\texttt{header}, \texttt{nav}, \texttt{main}, \texttt{section}, \texttt{footer}, \texttt{aside})
    \item Refondu votre \texttt{index.html} avec une structure sémantique propre
    \item Créé les \textbf{4 autres pages} de Jàng Afrig : \texttt{quiz.html}, \texttt{resultat.html}, \texttt{scores.html}, \texttt{apropos.html}
    \item Construit un \textbf{formulaire} (saisie du pseudo + sélection du mode)
    \item Lié les 5 pages entre elles avec une \textbf{navigation cohérente}
    \item \textbf{Validé} votre code HTML sur \href{https://validator.w3.org}{validator.w3.org}
    \item Posé le tag \textbf{\texttt{v0.2}} sur GitHub
    \item Rédigé votre rapport de Lab
\end{enumerate}
\end{objectifbox}

\tableofcontents

\newpage

% ===========================================================================
\section{Avant de coder}
% ===========================================================================

\subsection{Là où on en est}

Au Lab 1, vous avez créé une \textbf{page unique} \texttt{index.html} avec
toute l'information : titre, modes, palmarès, image, liens. Tout sur une
seule page.

\subsection{Le problème de la page unique}

Une vraie application a \textbf{plusieurs pages}. Imaginez Jàng Afrig
quand l'app sera complète :

\begin{itemize}[nosep]
    \item L'utilisateur arrive sur l'accueil
    \item Il choisit un mode et son pseudo
    \item Il clique «\,Commencer\,» $\to$ il va sur la page \emph{quiz}
    \item Il répond à 10 questions
    \item Il arrive sur la page \emph{résultat} avec son score
    \item Il peut aller voir la page \emph{mes scores} (classement)
    \item Il peut consulter la page \emph{à propos} (sources, équipe)
\end{itemize}

\textbf{5 pages distinctes}, donc \textbf{5 fichiers HTML} dans votre repo.

\subsection{Ce qu'on va construire aujourd'hui}

Les \textbf{5 pages HTML} de Jàng Afrig, sans CSS ni JavaScript. Toutes
liées entre elles avec une \textbf{barre de navigation} commune.

\begin{figure}[!ht]
\centering
\begin{tikzpicture}[
    page/.style={rounded corners=6pt, draw, thick, minimum width=2.6cm, minimum height=1.4cm, align=center, font=\small\bfseries},
    fleche/.style={-{Stealth[length=2.5mm]}, thick, line width=0.8pt, color=gray}
]
    \node[page, fill=pageaccueil!15, draw=pageaccueil] (accueil) {Accueil\\\normalfont\scriptsize\texttt{index.html}};
    \node[page, fill=pagequiz!15, draw=pagequiz, right=2cm of accueil] (quiz) {Quiz\\\normalfont\scriptsize\texttt{quiz.html}};
    \node[page, fill=pageresultat!15, draw=pageresultat, right=2cm of quiz] (resultat) {Résultat\\\normalfont\scriptsize\texttt{resultat.html}};
    \node[page, fill=pagescores!15, draw=pagescores, below=1.5cm of quiz] (scores) {Mes Scores\\\normalfont\scriptsize\texttt{scores.html}};
    \node[page, fill=pageapropos!15, draw=pageapropos, below=1.5cm of resultat] (apropos) {À propos\\\normalfont\scriptsize\texttt{apropos.html}};

    % Flux principal du quiz
    \draw[fleche, thick, color=pagequiz] (accueil) -- node[above, font=\scriptsize\itshape, fill=white, inner sep=1pt] {Commencer} (quiz);
    \draw[fleche, thick, color=pageresultat] (quiz) -- node[above, font=\scriptsize\itshape, fill=white, inner sep=1pt] {10 questions} (resultat);

    % Liens nav (toutes pages se lient mutuellement via la nav)
    \draw[fleche, dashed] (accueil.south) -- (scores.west);
    \draw[fleche, dashed] (accueil.south) -- (apropos.west);
    \draw[fleche, dashed] (resultat.south) -- (scores.east);
    \draw[fleche, dashed] (scores.east) -- (apropos.west);

    \node[below=0.3cm of scores, font=\scriptsize\itshape, text=gray] (lib) {(la nav présente sur toutes les pages permet d'aller partout)};
\end{tikzpicture}
\caption{Les 5 pages de Jàng Afrig et leur navigation.}
\label{fig:pages}
\end{figure}

\newpage

% ===========================================================================
\section{Les balises sémantiques HTML5}
% ===========================================================================

Au Lab 1, vous n'avez utilisé que des balises de \emph{contenu} :
\texttt{<h1>}, \texttt{<p>}, \texttt{<ul>}, etc. Mais HTML5 propose aussi
des balises de \textbf{structure}, dites \emph{sémantiques}.

\subsection{Pourquoi «\,sémantique\,» ?}

\textbf{Sémantique} signifie «\,qui a du sens\,». Au lieu de découper sa
page en boîtes anonymes (\texttt{<div>}, \texttt{<div>}, \texttt{<div>}),
on utilise des balises qui \textbf{disent ce que la zone contient} :
header, navigation, contenu principal, footer\ldots

\begin{infobox}[L'analogie du journal]
Pensez à un journal papier comme \emph{Le Soleil} ou \emph{L'Observateur} :
\begin{itemize}[nosep]
    \item En haut : la \textbf{une} (\texttt{<header>}) avec le logo et la date
    \item Sous la une : un \textbf{sommaire} (\texttt{<nav>}) qui dit où trouver quoi
    \item Au milieu : les \textbf{articles} (\texttt{<main>} contenant des \texttt{<article>})
    \item En bas : la \textbf{mention légale} (\texttt{<footer>}) avec contact et imprimeur
\end{itemize}
On pourrait tout mettre dans des boîtes anonymes, mais ce serait beaucoup
moins clair.
\end{infobox}

\subsection{Les 6 balises à connaître}

\begin{center}
\begin{tabular}{@{}lp{8cm}@{}}
\toprule
\textbf{Balise} & \textbf{Rôle} \\
\midrule
\texttt{<header>} & En-tête de la page (logo, titre principal, parfois la nav) \\
\addlinespace
\texttt{<nav>}    & Barre de navigation (les liens vers les autres pages) \\
\addlinespace
\texttt{<main>}   & Contenu principal de la page. \textbf{Une seule par page.} \\
\addlinespace
\texttt{<section>} & Une section thématique du contenu (\emph{ex.} «\,Modes de jeu\,») \\
\addlinespace
\texttt{<article>} & Un contenu \emph{autonome} qui pourrait être réutilisé ailleurs (un article de blog, une carte produit, un quiz\ldots) \\
\addlinespace
\texttt{<footer>}  & Pied de page (mentions, copyright, liens secondaires) \\
\addlinespace
\texttt{<aside>}   & Contenu lié mais pas essentiel (encart, pub, citation latérale) \\
\bottomrule
\end{tabular}
\end{center}

\subsection{Sémantique vs. anonyme : le même rendu, du sens en plus}

Visuellement, il n'y a aucune différence entre :

\begin{lstlisting}[language=HTML5]
<div>
    <div>Mon site</div>
    <div>Lien1 Lien2</div>
</div>
\end{lstlisting}

et :

\begin{lstlisting}[language=HTML5]
<header>
    <h1>Mon site</h1>
    <nav>Lien1 Lien2</nav>
</header>
\end{lstlisting}

Mais le second est compris par les moteurs de recherche, les lecteurs
d'écran, les outils d'analyse. \textbf{Un code sémantique est un code
respectueux de l'utilisateur et professionnel.}

\newpage

% ===========================================================================
\section{Étape 1 : Refondre index.html}
% ===========================================================================

On reprend votre \texttt{index.html} du Lab 1 et on \textbf{l'enveloppe}
dans une structure sémantique. Le contenu reste le même, on ajoute juste
des balises autour.

\subsection{Le nouveau squelette}

\textbf{Remplacez tout le \texttt{<body>}} de votre \texttt{index.html}
par cette structure :

\begin{lstlisting}[language=HTML5]
<body>

    <header>
        <h1>Jàng Afrig</h1>
        <p>Apprendre l'Afrique, un pays à la fois</p>
    </header>

    <nav>
        <a href="index.html">Accueil</a> |
        <a href="quiz.html">Jouer</a> |
        <a href="scores.html">Mes Scores</a> |
        <a href="apropos.html">À propos</a>
    </nav>

    <main>

        <section>
            <h2>Bienvenue</h2>
            <p>
                Bienvenue sur <strong>Jàng Afrig</strong>, un quiz géographique sur les
                <strong>54 pays africains</strong>. Choisissez un mode de jeu et battez
                votre meilleur score !
            </p>
        </section>

        <section>
            <h2>Les modes de jeu</h2>
            <ul>
                <li><strong>Capitales</strong> : « Quelle est la capitale du Mali ? »</li>
                <li><strong>Drapeaux</strong> : « À quel pays appartient ce drapeau ? »</li>
                <li><strong>Langues</strong> : « Dans quel pays parle-t-on le swahili ? »</li>
                <li><strong>Voisins</strong> : « Combien de voisins a l'Algérie ? »</li>
                <li><strong>Chrono</strong> : maximum de bonnes réponses en 60 secondes</li>
            </ul>
        </section>

        <section>
            <h2>Palmarès des champions</h2>
            <p><a href="scores.html">Voir le classement complet</a></p>
        </section>

    </main>

    <footer>
        <p>
            &copy; 2025--2026 Jàng Afrig --- Projet pédagogique L1 DUT1-Info, ESP/UCAD.
            <a href="apropos.html">En savoir plus</a>.
        </p>
    </footer>

</body>
\end{lstlisting}

\textbf{Enregistrez}. Visuellement la page est presque identique à celle
du Lab 1, mais elle est maintenant \textbf{structurée sémantiquement}. \cmark

\begin{retenirbox}[Ce qui a changé]
\begin{itemize}[nosep]
    \item Le bloc titre est dans un \texttt{<header>}
    \item Les liens vers les pages sont dans un \texttt{<nav>}
    \item Tout le contenu central est dans un \texttt{<main>}, divisé en 3 \texttt{<section>}
    \item Les mentions sont dans un \texttt{<footer>}
    \item Le tableau du palmarès est remplacé par un lien vers \texttt{scores.html} (qu'on va créer)
\end{itemize}
\end{retenirbox}

\subsection{L'entité \texttt{\&copy;}}

Vous avez vu \texttt{\&copy;} dans le footer. C'est une \textbf{entité HTML} :
elle s'affiche \texttt{©}. Quelques entités utiles :

\begin{center}
\begin{tabular}{@{}lll@{}}
\toprule
\textbf{Entité} & \textbf{Affiche} & \textbf{Quand l'utiliser} \\
\midrule
\texttt{\&copy;}  & ©       & Copyright \\
\texttt{\&nbsp;}  & espace  & Espace insécable (avant : ; ! ?) \\
\texttt{\&amp;}   & \&      & Le caractère «\,et commercial\,» \\
\texttt{\&lt;}    & <       & Le chevron ouvrant (sans démarrer une balise) \\
\texttt{\&gt;}    & >       & Le chevron fermant \\
\texttt{\&hellip;} & \ldots & Points de suspension \\
\bottomrule
\end{tabular}
\end{center}

\newpage

% ===========================================================================
\section{Étape 2 : La page Quiz avec un formulaire}
% ===========================================================================

Créez un nouveau fichier \texttt{quiz.html} dans votre repo. C'est la page
où l'utilisateur saisit son pseudo et choisit son mode de jeu.

\subsection{Code complet de quiz.html}

\begin{lstlisting}[language=HTML5]
<!DOCTYPE html>
<html lang="fr">
    <head>
        <meta charset="UTF-8">
        <title>Quiz --- Jàng Afrig</title>
    </head>
    <body>

        <header>
            <h1>Jàng Afrig</h1>
            <p>Apprendre l'Afrique, un pays à la fois</p>
        </header>

        <nav>
            <a href="index.html">Accueil</a> |
            <a href="quiz.html">Jouer</a> |
            <a href="scores.html">Mes Scores</a> |
            <a href="apropos.html">À propos</a>
        </nav>

        <main>
            <section>
                <h2>Avant de jouer</h2>

                <form>
                    <p>
                        <label for="pseudo">Votre pseudo :</label>
                        <input type="text" id="pseudo" name="pseudo"
                               placeholder="Ex. Fatou" required>
                    </p>

                    <p>
                        <label for="mode">Mode de jeu :</label>
                        <select id="mode" name="mode" required>
                            <option value="">-- Choisir un mode --</option>
                            <option value="capitales">Capitales</option>
                            <option value="drapeaux">Drapeaux</option>
                            <option value="langues">Langues</option>
                            <option value="voisins">Voisins</option>
                            <option value="chrono">Chrono</option>
                        </select>
                    </p>

                    <p>
                        <button type="submit">Commencer le quiz</button>
                    </p>
                </form>
            </section>
        </main>

        <footer>
            <p>&copy; 2025--2026 Jàng Afrig --- L1 DUT1-Info, ESP/UCAD.</p>
        </footer>

    </body>
</html>
\end{lstlisting}

\textbf{Enregistrez}. Live Server affiche la page (cliquez sur «\,Jouer\,» depuis
l'accueil). Vous voyez un formulaire avec un champ texte, un menu
déroulant, et un bouton. \cmark

\subsection{Décortiquer le formulaire}

\begin{center}
\begin{tabular}{@{}lp{8.5cm}@{}}
\toprule
\textbf{Élément} & \textbf{Rôle} \\
\midrule
\texttt{<form>} & Conteneur du formulaire (englobe tous les champs) \\
\addlinespace
\texttt{<label for="X">} & Étiquette d'un champ. \texttt{for="X"} se lie à \texttt{id="X"}. \\
\addlinespace
\texttt{<input type="text">} & Champ de texte sur une ligne \\
\addlinespace
\texttt{placeholder="..."} & Texte gris d'aide qui disparaît à la frappe \\
\addlinespace
\texttt{required}     & Le champ est obligatoire (validation HTML5 native) \\
\addlinespace
\texttt{<select>} + \texttt{<option>} & Menu déroulant à options fixes \\
\addlinespace
\texttt{<button type="submit">} & Bouton qui soumet le formulaire \\
\bottomrule
\end{tabular}
\end{center}

\begin{infobox}[Test du formulaire]
Cliquez sur «\,Commencer le quiz\,» sans rien remplir : le navigateur
\textbf{bloque} l'envoi et affiche un message «\,Veuillez remplir ce
champ\,». C'est la \emph{validation HTML5 native}, gratuite, sans
JavaScript.
\end{infobox}

\begin{attentionbox}[Le formulaire ne mène nulle part pour l'instant]
Cliquer sur «\,Commencer\,» va recharger la page. C'est normal : il
manque le JavaScript pour traiter les données. On s'en occupera au
\textbf{Lab 7}. Pour aujourd'hui, l'objectif est juste d'avoir le
formulaire \emph{visible et structuré}.
\end{attentionbox}

\subsection{Pourquoi \texttt{<label>} est important}

L'attribut \texttt{for} d'un \texttt{<label>} relie l'étiquette au champ.
Cliquer sur «\,Votre pseudo :\,» met le curseur dans le champ. Pour les
utilisateurs malvoyants, le lecteur d'écran annonce «\,Votre pseudo,
champ texte\,» au lieu d'un anonyme «\,champ texte\,». \textbf{Toujours
mettre un \texttt{<label>}.}

\newpage

% ===========================================================================
\section{Étape 3 : Les 3 autres pages}
% ===========================================================================

On va créer rapidement les 3 pages restantes. Elles ont toutes la \textbf{même
structure} (header + nav + main + footer), seul le contenu du \texttt{<main>}
change.

\subsection{Page resultat.html}

Créez \texttt{resultat.html} avec ce contenu (la partie \texttt{<main>}
seulement est différente, le reste est identique aux autres pages) :

\begin{lstlisting}[language=HTML5]
<main>
    <section>
        <h2>Résultat de votre partie</h2>

        <p>Bravo <strong>Fatou</strong> !</p>
        <p>Mode : <em>Capitales</em></p>
        <p>Score : <strong>8 / 10</strong></p>

        <p>
            <a href="quiz.html">Rejouer</a> |
            <a href="scores.html">Voir mes scores</a> |
            <a href="index.html">Retour à l'accueil</a>
        </p>
    </section>
</main>
\end{lstlisting}

\textbf{Note :} les valeurs «\,Fatou\,», «\,Capitales\,», «\,8/10\,» sont
\emph{hardcodées} pour l'instant. Au Lab 7, JavaScript les remplacera par
les vraies valeurs de la partie.

\subsection{Page scores.html}

Créez \texttt{scores.html}. Le \texttt{<main>} contient le tableau du
palmarès (récupéré du Lab 1) :

\begin{lstlisting}[language=HTML5]
<main>
    <section>
        <h2>Palmarès des champions</h2>

        <p>Les meilleurs scores enregistrés (données fictives pour l'instant) :</p>

        <table>
            <tr>
                <th>Rang</th>
                <th>Pseudo</th>
                <th>Mode</th>
                <th>Score</th>
                <th>Date</th>
            </tr>
            <tr>
                <td>1</td>
                <td>Fatou</td>
                <td>Capitales</td>
                <td>10/10</td>
                <td>15 mars 2026</td>
            </tr>
            <tr>
                <td>2</td>
                <td>Moussa</td>
                <td>Drapeaux</td>
                <td>9/10</td>
                <td>12 mars 2026</td>
            </tr>
            <tr>
                <td>3</td>
                <td>Aminata</td>
                <td>Capitales</td>
                <td>9/10</td>
                <td>10 mars 2026</td>
            </tr>
            <tr>
                <td>4</td>
                <td>Cheikh</td>
                <td>Voisins</td>
                <td>8/10</td>
                <td>8 mars 2026</td>
            </tr>
            <tr>
                <td>5</td>
                <td>Awa</td>
                <td>Langues</td>
                <td>8/10</td>
                <td>5 mars 2026</td>
            </tr>
        </table>
    </section>
</main>
\end{lstlisting}

\subsection{Page apropos.html}

Créez \texttt{apropos.html}. Le \texttt{<main>} :

\begin{lstlisting}[language=HTML5]
<main>
    <section>
        <h2>À propos de Jàng Afrig</h2>
        <p>
            Jàng Afrig est une application web de quiz géographique sur les
            <strong>54 pays africains</strong>. Le mot <em>jàng</em> signifie
            « apprendre, étudier » en wolof.
        </p>
        <p>
            Ce projet est développé dans le cadre du cours
            <strong>Fondamentaux du Web</strong> en L1 DUT1-Informatique
            à l'École Supérieure Polytechnique (ESP/UCAD).
        </p>
    </section>

    <section>
        <h2>Sources des données</h2>
        <p>
            Les données sur les pays sont fournies par
            <a href="https://restcountries.com">REST Countries v3.1</a>,
            une API libre et gratuite.
        </p>
    </section>

    <section>
        <h2>Technologies utilisées</h2>
        <ul>
            <li>HTML5 pour la structure</li>
            <li>CSS3 pour la mise en forme</li>
            <li>JavaScript moderne pour l'interactivité</li>
            <li>Git et GitHub pour le versioning et l'hébergement</li>
        </ul>
    </section>

    <section>
        <h2>Code source</h2>
        <p>
            Le code est disponible sur GitHub :
            <a href="https://github.com/votre-username/jang-afrig">
                github.com/votre-username/jang-afrig
            </a>
        </p>
    </section>
</main>
\end{lstlisting}

\begin{retenirbox}[Astuce de productivité]
Au lieu de retaper le squelette à chaque fois, \textbf{copiez \texttt{quiz.html}}
puis \textbf{renommez} la copie en \texttt{resultat.html}, et changez
seulement le contenu du \texttt{<main>}. Pareil pour \texttt{scores.html}
et \texttt{apropos.html}. C'est ainsi qu'on travaille en vrai.
\end{retenirbox}

\newpage

% ===========================================================================
\section{Étape 4 : Tester la navigation}
% ===========================================================================

\subsection{Recharger la page d'accueil}

Dans le navigateur (où Live Server tourne) :
\begin{enumerate}[nosep]
    \item Allez sur \url{http://127.0.0.1:5500/index.html}
    \item Cliquez sur «\,Jouer\,» dans la nav $\to$ vous arrivez sur \texttt{quiz.html}
    \item Cliquez sur «\,Mes Scores\,» $\to$ \texttt{scores.html}
    \item Cliquez sur «\,À propos\,» $\to$ \texttt{apropos.html}
    \item Depuis n'importe quelle page, cliquez sur «\,Accueil\,» $\to$ retour à \texttt{index.html}
\end{enumerate}

\textbf{Toutes les pages sont accessibles depuis n'importe quelle autre.} \cmark

\subsection{Vérification : la nav doit être identique partout}

\begin{attentionbox}[Le piège du copier-coller]
Si vous changez le menu sur une page, vous devez le changer sur toutes
les autres. Sinon, l'utilisateur a une nav différente selon où il est :
incohérent et frustrant.

\medskip

\textbf{Aujourd'hui}, on accepte cette duplication (4 fois le même
\texttt{<nav>}). \textbf{Plus tard} (en L2 avec un framework, ou même au
Lab 7 avec JS), on apprendra à ne définir la nav qu'\emph{une seule fois}.
\end{attentionbox}

\newpage

% ===========================================================================
\section{Valider votre code HTML}
% ===========================================================================

Le W3C (World Wide Web Consortium, l'organisme qui définit HTML) propose
un \textbf{validateur en ligne gratuit} qui détecte les erreurs de syntaxe
dans votre HTML.

\subsection{Comment l'utiliser}

\begin{enumerate}[nosep]
    \item Allez sur \url{https://validator.w3.org/nu/}
    \item Onglet \textbf{Validate by Direct Input}
    \item Ouvrez \texttt{index.html} dans VS Code, sélectionnez \textbf{tout} (\texttt{Ctrl+A}), copiez (\texttt{Ctrl+C})
    \item Collez dans la zone de texte du validateur
    \item Cliquez \textbf{Check}
\end{enumerate}

\subsection{Lire le rapport}

\begin{itemize}[leftmargin=*,nosep]
    \item Bandeau \textbf{vert} : «\,Document checking completed. No errors or warnings to show.\,» $\to$ votre HTML est parfait \cmark
    \item Bandeau \textbf{rouge} : il y a des erreurs. Le rapport indique la ligne et la nature du problème.
\end{itemize}

\textbf{Faites tourner le validateur sur les 5 pages.} Corrigez les
erreurs au fur et à mesure.

\begin{infobox}[Erreurs courantes]
\begin{itemize}[leftmargin=*,nosep]
    \item Balise non fermée : \texttt{<p>texte} sans \texttt{</p>}
    \item Balise mal imbriquée : \texttt{<p><strong>texte</p></strong>}
    \item Attribut \texttt{alt} manquant sur \texttt{<img>}
    \item Plusieurs \texttt{<h1>} dans la même page
    \item Caractères spéciaux non échappés (\texttt{<} au lieu de \texttt{\&lt;} dans du texte)
\end{itemize}
\end{infobox}

\newpage

% ===========================================================================
\section{Exercices d'application}
% ===========================================================================

\begin{objectifbox}[Trois mini-exercices]
\textbf{Exercice 1 --- Ajouter un \texttt{<aside>} sur l'accueil}

Sur \texttt{index.html}, ajoutez un \texttt{<aside>} (encart latéral) après
le \texttt{<main>} avec un «\,saviez-vous ?\,» géographique. Exemple :

\begin{lstlisting}[language=HTML5]
<aside>
    <h3>Le saviez-vous ?</h3>
    <p>
        L'Afrique compte <strong>54 pays</strong> reconnus par l'ONU,
        avec plus de <strong>2000 langues</strong> parlées.
    </p>
</aside>
\end{lstlisting}

\medskip

\textbf{Exercice 2 --- Enrichir le formulaire}

Sur \texttt{quiz.html}, ajoutez dans le \texttt{<form>} :
\begin{itemize}[nosep]
    \item Un champ \texttt{<input type="number">} pour choisir le \textbf{nombre de questions} (min 5, max 20, valeur par défaut 10)
    \item 3 \texttt{<input type="radio">} pour la \textbf{difficulté} : Facile, Moyen, Difficile
    \item N'oubliez pas les \texttt{<label>} associés
\end{itemize}

\medskip

\textbf{Exercice 3 --- Page formulaire de contact (bonus)}

Créez une nouvelle page \texttt{contact.html} avec un formulaire complet :
\begin{itemize}[nosep]
    \item Nom (\texttt{input type="text"}, required)
    \item Email (\texttt{input type="email"}, required)
    \item Sujet (\texttt{select} avec 3 options : Question, Bug, Suggestion)
    \item Message (\texttt{<textarea>} de 5 lignes, required)
    \item Bouton «\,Envoyer\,»
\end{itemize}
Ajoutez le lien vers cette page dans le \texttt{<nav>} de toutes les pages.

\medskip

\textbf{Validez vos pages au W3C} avant de passer à la suite.
\end{objectifbox}

\newpage

% ===========================================================================
\section{Aide-mémoire}
% ===========================================================================

\subsection{Les balises sémantiques}

\begin{center}
\begin{tabular}{@{}lp{7.5cm}@{}}
\toprule
\textbf{Balise}     & \textbf{Rôle} \\
\midrule
\texttt{<header>}   & En-tête (logo, titre principal) \\
\texttt{<nav>}      & Barre de navigation \\
\texttt{<main>}     & Contenu principal (un seul par page) \\
\texttt{<section>}  & Section thématique \\
\texttt{<article>}  & Contenu autonome réutilisable \\
\texttt{<aside>}    & Contenu lié mais secondaire \\
\texttt{<footer>}   & Pied de page \\
\bottomrule
\end{tabular}
\end{center}

\subsection{Les balises de formulaire}

\begin{center}
\begin{tabular}{@{}lp{7.5cm}@{}}
\toprule
\textbf{Balise}                           & \textbf{Rôle} \\
\midrule
\texttt{<form>}                           & Conteneur d'un formulaire \\
\texttt{<label for="X">}                  & Étiquette d'un champ \\
\texttt{<input type="text">}              & Champ texte \\
\texttt{<input type="email">}             & Champ email (validation auto) \\
\texttt{<input type="number">}            & Champ numérique \\
\texttt{<input type="radio">}             & Bouton radio (choix unique) \\
\texttt{<input type="checkbox">}          & Case à cocher (choix multiple) \\
\texttt{<select>} + \texttt{<option>}     & Menu déroulant \\
\texttt{<textarea>}                       & Champ texte multi-lignes \\
\texttt{<button type="submit">}           & Bouton de soumission \\
\bottomrule
\end{tabular}
\end{center}

\subsection{Attributs utiles}

\begin{center}
\begin{tabular}{@{}lp{7.5cm}@{}}
\toprule
\textbf{Attribut}         & \textbf{Rôle} \\
\midrule
\texttt{id="..."}         & Identifiant unique sur la page (référencé par \texttt{label for}) \\
\texttt{name="..."}       & Nom du champ (utilisé côté serveur, plus tard) \\
\texttt{placeholder="..."} & Texte d'aide gris dans un champ vide \\
\texttt{required}         & Champ obligatoire \\
\texttt{value="..."}      & Valeur par défaut d'un champ \\
\bottomrule
\end{tabular}
\end{center}

\newpage

% ===========================================================================
\section{Routine Git de fin de Lab}
% ===========================================================================

Dans le terminal intégré de VS Code :

\begin{lstlisting}
git status
git add .
git commit -m "v0.2 : 5 pages HTML semantiques + formulaire"
git push
git tag v0.2
git push --tags
\end{lstlisting}

Vérifiez sur \url{https://github.com/votre-username/jang-afrig} que :
\begin{itemize}[nosep]
    \item Les 5 fichiers HTML sont là (\texttt{index.html}, \texttt{quiz.html}, \texttt{resultat.html}, \texttt{scores.html}, \texttt{apropos.html})
    \item Le tag \texttt{v0.2} est visible dans l'onglet \textbf{Tags}
\end{itemize}

% ===========================================================================
\section{Rapport à produire}
% ===========================================================================

\begin{attentionbox}[Rapport du Lab 2 --- obligatoire]
\textbf{Format :} 1-2 pages, déposé dans \texttt{rapports/rapport-lab-2-NOM-Prenom.pdf}
de votre repo.

\medskip

\textbf{Contenu (4 sections) :}
\begin{enumerate}[nosep]
    \item \textbf{Ce que j'ai fait} : 5 pages créées, balises sémantiques utilisées, formulaire construit, exercices d'application complétés, validation W3C passée.
    \item \textbf{Difficultés rencontrées} : navigation cassée, validateur W3C qui rouspète, formulaire qui se comporte bizarrement\ldots
    \item \textbf{Comment j'ai résolu ces difficultés} : recherche, aide d'un camarade, lecture de la doc.
    \item \textbf{Lien vers mon repo GitHub avec le tag v0.2} : \\
          ex. \url{https://github.com/fatou-diop/jang-afrig/releases/tag/v0.2}
\end{enumerate}
\end{attentionbox}

% ===========================================================================
\section{Résumé et transition}
% ===========================================================================

\subsection{Ce que vous avez accompli}

\begin{itemize}[leftmargin=*,nosep]
    \item Maîtrisé les balises sémantiques HTML5 \cmark
    \item Construit les 5 pages de Jàng Afrig \cmark
    \item Construit votre premier formulaire fonctionnel \cmark
    \item Lié les pages avec une navigation cohérente \cmark
    \item Validé votre code au W3C \cmark
    \item Posé le tag \texttt{v0.2} \cmark
\end{itemize}

\subsection{Le bloc HTML est terminé}

Vous avez vu \textbf{tout ce qu'il faut savoir en HTML} pour ce semestre :
structure, contenu, sémantique, formulaires, tableaux, liens, images.

\subsection{Ce qu'on fera au CM2 et au Lab 3}

\textbf{CM2 (S4) :} le CSS arrive. Vous découvrirez les sélecteurs, la
cascade, le box model, et Flexbox.

\textbf{Lab 3 (S5) :} vous habillerez Jàng Afrig avec des couleurs
panafricaines, des polices Google Fonts, et de jolis boutons. Tag \texttt{v0.3}.

\begin{projetbox}[À ce stade du semestre]
Votre application a une structure complète et propre. Elle est encore
moche (texte noir sur fond blanc), mais \textbf{toutes les fondations
HTML sont posées}. Aux Labs 3 et 4, on transforme ça en application
visuellement attractive.
\end{projetbox}

\vspace{0.5cm}

\begin{center}
\fbox{
\begin{minipage}{0.85\textwidth}
\centering
\vspace{0.5cm}
\textbf{\Large Lab 2 --- Terminé \cmark}\\[0.3cm]
\normalsize
Vous avez construit l'ossature complète de Jàng Afrig.\\
Place à la décoration au CM2 et au Lab 3.\\[0.2cm]
\textit{Rendez-vous au CM2 !}\\[0.3cm]
\vspace{0.3cm}
\end{minipage}
}
\end{center}

\end{document}
